\section{Final checklist}

This chapter completed the transition from probability models to statistical
inference. The essential chain is
\[
\text{model}
\longrightarrow
\text{sampling distribution}
\longrightarrow
\text{large-sample regularity}
\longrightarrow
\text{observed statistic}
\longrightarrow
\text{statistical conclusion}.
\]

The conclusion is never produced by the observed number alone. It depends on the
sampling mechanism, the assumptions, and the probability law used to describe
how the statistic would vary across possible samples.

\subsection*{Conceptual understanding}

You should be able to explain the following distinctions:

\begin{itemize}
    \item probability reasoning versus statistical inference;
    \item a model parameter \(\theta\) versus an estimator \(\widehat\theta\);
    \item convergence of a statistic versus exact equality with the parameter;
    \item the law of large numbers versus the central limit theorem;
    \item population standard deviation versus standard error;
    \item a theoretical standard error versus an estimated standard error;
    \item a confidence level versus a probability assigned to a fixed parameter;
    \item failure to reject a hypothesis versus proof that it is true;
    \item statistical significance versus practical importance.
\end{itemize}

\subsection*{The main regularities}

For independent identically distributed observations with finite mean \(\mu\),
the law of large numbers gives
\[
\overline X_n\xrightarrow{P}\mu.
\]
It explains why large-sample averages become stable around the population mean.

If the observations also have finite positive variance \(\sigma^2\), the central
limit theorem gives
\[
\frac{\overline X_n-\mu}{\sigma/\sqrt n}
\xrightarrow{d}N(0,1).
\]
It describes the scale and approximate shape of the remaining fluctuations.

The two results have different roles:
\[
\text{LLN: where the statistic stabilizes,}
\]
\[
\text{CLT: how the statistic fluctuates around that value.}
\]

\subsection*{Standard errors}

The standard error of a statistic is the standard deviation of its sampling
distribution:
\[
\operatorname{SE}(T)=\sqrt{\Var(T)}.
\]

For a sample mean,
\[
\operatorname{SE}(\overline X)=\frac{\sigma}{\sqrt n}.
\]

For a sample proportion,
\[
\operatorname{SE}(\widehat p)
=
\sqrt{\frac{p(1-p)}{n}}.
\]

Standard errors measure uncertainty caused by random sampling. They do not
measure all possible errors in data collection or modelling.

\subsection*{Confidence intervals}

A large-sample confidence interval has the general form
\[
\text{estimate}
\pm
\text{critical value}\times\text{estimated standard error}.
\]

For an approximate \(95\%\) interval, the normal critical value is often
\[
1.96.
\]
The confidence level describes the long-run coverage of the interval-producing
procedure under repeated sampling.

A confidence interval can also be viewed as a range of parameter values reasonably
compatible with the observed statistic and the assumed sampling model.

\subsection*{Hypothesis tests}

A null hypothesis has the form
\[
H_0:\theta=\theta_0.
\]
The hypothesis determines a sampling distribution for a test statistic. The
observed statistic is then compared with this null distribution.

A standardized statistic often has the form
\[
Z
=
\frac{T-\theta_0}{\operatorname{SE}_{H_0}(T)}.
\]

The p-value is the probability, under \(H_0\), of obtaining a result at least as
extreme as the observed one. A small p-value indicates incompatibility between
the data and the null model.

Remember:
\[
p\text{-value}
\neq
P(H_0\text{ is true}\mid\text{data}).
\]

\subsection*{Typical tasks}

After studying this chapter, you should be able to:

\begin{itemize}
    \item explain why stable large-sample behaviour makes inference possible;
    \item distinguish the conclusions of the law of large numbers and the central
    limit theorem;
    \item standardize a sample mean or sample proportion;
    \item interpret a difference in standard-error units;
    \item calculate an approximate standard error for a mean or proportion;
    \item construct and interpret a basic large-sample confidence interval;
    \item formulate a simple null hypothesis;
    \item construct a standardized test statistic under the null model;
    \item interpret a p-value as a probability about hypothetical data;
    \item distinguish evidence against a model from logical proof;
    \item identify assumptions on which an inferential conclusion depends.
\end{itemize}

\subsection*{Common mistakes}

Avoid the following mistakes:

\begin{itemize}
    \item saying that the law of large numbers makes an average exactly equal to
    the parameter;
    \item saying that the central limit theorem makes the population normal;
    \item using a normal approximation without considering sample size or model
    assumptions;
    \item confusing the standard deviation of observations with the standard
    error of a statistic;
    \item interpreting a \(95\%\) confidence interval as assigning probability
    \(0.95\) to a fixed parameter after the interval is observed;
    \item interpreting a p-value as the probability that the null hypothesis is
    true;
    \item treating failure to reject as proof of no difference;
    \item treating statistical significance as evidence of practical importance;
    \item believing that a large sample removes bias, dependence, measurement
    error, or poor modelling.
\end{itemize}

\subsection*{Questions for self-explanation}

A good test of understanding is whether you can answer these questions clearly:

\begin{enumerate}
    \item Why does statistical inference reverse the direction of probability
    reasoning?
    \item Why can one observed statistic contain information about an unknown
    parameter?
    \item What does the law of large numbers say about sample averages?
    \item What additional information is supplied by the central limit theorem?
    \item Why does \(1/\sqrt n\) appear in standard errors?
    \item What is random in a confidence-interval procedure?
    \item What does a \(95\%\) confidence level mean under repeated sampling?
    \item How does a null hypothesis determine a sampling distribution?
    \item Why is an observed result several standard errors from the null centre
    evidence against the null model?
    \item What does a p-value measure, and what does it not measure?
    \item Why can a very large sample make an unimportant difference statistically
    significant?
    \item How can a failure of modelling assumptions affect a confidence interval
    or hypothesis test?
\end{enumerate}

\begin{summarybox}
Statistical inference is possible because statistics have predictable sampling
behaviour. The law of large numbers produces stability, the central limit theorem
describes fluctuations, and standard errors express their scale. Confidence
intervals and hypothesis tests use these regularities to connect one observed
sample with claims about an unknown model while keeping uncertainty explicit.
\end{summarybox}
